% TOP_PROBLEM: {"id": 3190, "title": "Linial-Berge path partition duality", "category_id": 3, "category": {"id": 3, "name": "graph_theory", "display_name": "Graph Theory", "description": "Problems involving graphs, networks, and their properties.", "slug": "graph-theory", "order_index": 3, "created_at": "2026-07-31T15:26:25.671Z"}, "status": "open", "problem_number": "OPG-170", "set_id": 12}
% FIRST_POSTED: 2026-10-01
% ENTRY_KIND: statement_only
% UnsolvedMath ID 3190; source catalogue code OPG-170
% !TeX program = lualatex
% Definition and sources only; this file is not an LLM solution attempt.
\documentclass[11pt,a4paper]{article}
\usepackage[margin=25mm]{geometry}
\usepackage{fontspec}
\setmainfont{lmroman10-regular.otf}[BoldFont=lmroman10-bold.otf,
 ItalicFont=lmroman10-italic.otf,BoldItalicFont=lmroman10-bolditalic.otf]
\usepackage{amsmath,amssymb,mathtools}
\usepackage{unicode-math}
\setmathfont{latinmodern-math.otf}
\AtBeginDocument{\let\setminus\smallsetminus}
\usepackage{xurl,hyperref}
\hypersetup{colorlinks=true,urlcolor=blue}
\setcounter{secnumdepth}{0}
\setlength{\parindent}{0pt}
\setlength{\parskip}{5pt}
\setlength{\emergencystretch}{3em}
\begin{document}
\section*{Linial-Berge path partition duality}\label{3190}
{\small UnsolvedMath ID 3190 · OPG-170 · Graph Theory · OpenGarden}
\subsection{Definitions and mathematical statement}
Let $D=(V,A)$ be a finite loopless digraph, allowing both opposite arcs between two vertices. A directed path is a sequence of distinct vertices with an arc from each vertex to its successor. A singleton vertex is also a path. A path partition $\mathcal P$ is a family of vertex-disjoint directed paths whose vertex sets partition $V$.

For an integer $k\ge1$, define the $k$-norm of a partition and its minimum by
\[
 \|\mathcal P\|_k=\sum_{P\in\mathcal P}\min\{|V(P)|,k\},
 \qquad \pi_k(D)=\min_{\mathcal P}\|\mathcal P\|_k.
\]
An independent set in $D$ is a vertex set with no arc between any two of its distinct vertices, in either direction. Let $\alpha_k(D)$ be the maximum size of a union of $k$ independent sets. Equivalently, it is the largest number of vertices inducing a graph whose underlying undirected graph has a proper coloring with at most $k$ colors. Empty color classes are permitted.

The conjecture is $\pi_k(D)\le\alpha_k(D)$ for every finite $D$ and every positive integer $k$. Thus some partition should have total truncated path length at most the largest $k$-colorable vertex set. The independent sets do not merely induce acyclic digraphs: arcs inside one color class are prohibited altogether. Path lengths in the formula count vertices, not arcs, so singleton paths contribute one. Both optimization quantities are finite and attained.
\subsection{Short English statement}
Can every digraph be partitioned into directed paths whose total k-truncated vertex counts are bounded by its largest union of k independent sets?
\subsection{Sources}
\begin{itemize}
\item[S1] ULAM.AI and the UnsolvedMath Contributors, supplied problems.json, record 3190 (OPG-170), OpenGarden. Catalogue identity and source transcription; CC BY 4.0.
\url{https://huggingface.co/datasets/ulamai/UnsolvedMath}.
\item[S2] Open Problem Garden, Linial-Berge path partition duality. Original problem and discussion.
\url{http://www.openproblemgarden.org/op/linial_berge_path_partition_duality}.
\end{itemize}
\end{document}
